% Ad Familiares 11.10 (ad-familiares-11-10) --- English translation
% Translated by Claude (Anthropic) from the Latin (Perseus).
% Style guide: STYLE.md.

\ciceroLetterOpener{D. BRVTVS S. D. M. CICERONI}

\ciceroSection{1} I do not think the commonwealth owes more to me than I owe to you. That I can be more grateful to you than those perverse fellows are unjust to me, you may take as established --- though whether this is even a fit moment to be saying so, I leave aside --- and that I would rather have your judgement than the judgements of all that other lot on the opposite side; for you judge of us from a sure and true sense, while malice and envy keep them from doing the same in the highest degree. Let them obstruct me as much as they like in the way of honours, provided they do not obstruct the commonwealth from being properly administered by me. In what danger that stands, I shall set forth as briefly as I can.

\ciceroSection{2} First of all, how great a disturbance in the affairs of the City the death of the consuls brings on, and how great an appetite the vacancy puts into men's minds, does not escape you. I think I have written enough on what can be entrusted to a letter; for I know whom I am writing to.

\ciceroSection{3} I return now to Antony. Who, when he came out of the rout with a small band of unarmed foot, by emptying the slave-pens and snatching up men of every sort seems to have got together a fair-sized number. To this has been added the troop of Ventidius, which, having made its way across the Apennines by the most difficult of marches, has reached Vada, and there has joined itself to Antony. With Ventidius is a body of veterans and armed men of fair numbers.

\ciceroSection{4} Antony's plans must of necessity be these: either to take himself to Lepidus, if he is received; or to hold himself in the Apennines and the Alps, and by raids with the cavalry, of which he has many, to lay waste the country into which he runs; or again to take himself back into Etruria, that part of Italy being without an army. But if Caesar had listened to me and crossed the Apennines, I should have driven Antony into such narrow straits that he would have been finished off by want rather than by the sword. But Caesar cannot be commanded, nor Caesar his own army --- both of which are very bad indeed. Such being the state of things, as for what concerns me personally, I do not, as I said above, stand in the way of men's obstructing. How these matters may be untangled --- or, when they shall have been untangled by you, may not then be tangled again --- is what I fear.

\ciceroSection{5} I can no longer feed my soldiers. When I entered upon the freeing of the commonwealth, I had four hundred thousand sesterces and more of my own. So far is it from being the case that anything of my own property is unencumbered, that I have by now bound all my friends in debt. Seven legions in number I am now feeding --- with what difficulty, you may judge for yourself; not if I had Varro's treasuries could I keep up with the expense. As soon as I have anything sure about Antony, I shall let you know. You will love me on this condition, that you find me doing the same by you. 5 May, from camp at Dertona.
